\documentclass[12pt]{article}
\usepackage{setspace}
\onehalfspacing

\usepackage{iftex}
\ifPDFTeX
  \usepackage[utf8]{inputenc}
  \usepackage[T1]{fontenc}
\else
  \usepackage{fontspec}
\fi
\usepackage[spanish, es-tabla]{babel}
\usepackage[margin=1in]{geometry}
\usepackage{enumitem}
\usepackage{hyperref}
\usepackage{microtype}
\usepackage{amsmath}
\usepackage{graphicx}
\usepackage{biblatex}
\addbibresource{refs.bib}

\hypersetup{
  colorlinks=true,
  linkcolor=black,
  urlcolor=blue,
  citecolor=black
}

\setlist[itemize]{leftmargin=*, itemsep=0.35em}
\newcommand{\institucion}{Universidad de Costa Rica}
\newcommand{\campus}{Ciudad Universitaria Rodrigo Facio}
\newcommand{\facultad}{Facultad de Ciencias Básicas}
\newcommand{\escuela}{Escuela de Física}
\newcommand{\curso}{Mecánica Estadística Computacional}
\newcommand{\tituloanteproyecto}{Proyecto de Investigación}
\newcommand{\titulo}{Simulación Monte Carlo de Casos Nuevos de COVID-19 por Zona Geográfica}
\newcommand{\autor}{Gustavo A. Calvo Gutiérrez}
\newcommand{\lugarfecha}{San José, 2026.}

\begin{document}

\begin{titlepage}
  \begin{center}
    \LARGE \textbf{\textsc{\institucion}}
    \Large
    \\[1cm]
    \textsc{\campus}
    \\[0.1cm]
    \textsc{\facultad}
    \\[0.1cm]
    \textsc{\escuela} \\
    \vspace*{20 mm}
    \LARGE{\textsc{\textbf{\curso}}}
  \end{center}

  \vspace{2cm}

  \hrulefill
  \begin{center}
    \large \textbf{\textsc{\tituloanteproyecto}}
    \\[0.5cm]
    \textsc{\titulo}
  \end{center}
  \hrulefill

  \vfill

  \begin{center}
    \large \textbf{\autor}
  \end{center}

  \vspace{1cm}

  \begin{center}
    \textit{\lugarfecha}
  \end{center}
\end{titlepage}

\section{Pregunta de Investigación y Justificación}

\subsection{Pregunta de Investigación}

\noindent
\textit{
  ¿Cómo puede un modelo de simulación Monte Carlo, aplicado por zona
  geográfica y basado en el procedimiento propuesto por Xie \cite{Xie2020},
  estimar la cantidad esperada de casos nuevos de COVID-19 a partir de series
  temporales observadas?
}

\subsection{Justificación}

El estudio de la propagación de COVID-19 continúa siendo relevante como caso de
referencia para comprender fenómenos epidémicos donde la transmisión depende de
interacciones sociales y de variaciones temporales en la dinámica de contagio.
Los modelos estocásticos ofrecen una ventaja importante frente a aproximaciones
puramente deterministas, ya que permiten representar explícitamente la
incertidumbre asociada a la ocurrencia de nuevos contagios y a la variabilidad
temporal de la transmisión.

El trabajo de Xie \cite{Xie2020} plantea una estrategia basada en simulación
Monte Carlo para describir la dinámica de contagio mediante tasas de infección
variables en el tiempo, junto con distribuciones probabilísticas para modelar
tanto la cantidad de personas contagiadas por cada caso activo como el momento
en que dichos contagios ocurren.

El presente anteproyecto propone tomar ese artículo como punto de partida
metodológico para realizar simulaciones por zona geográfica y estimar, a partir
de múltiples corridas aleatorias, la cantidad esperada de casos nuevos en cada
periodo. La meta no es evaluar la efectividad de medidas sanitarias ni comparar
intervenciones, sino construir una herramienta de simulación capaz de generar
trayectorias plausibles de contagio y resumirlas mediante su valor esperado.
Este enfoque es pertinente porque permite vincular la simulación con series
epidemiológicas observadas y estudiar la variabilidad inherente al proceso de
transmisión. Para ello, el estudio puede apoyarse en fuentes comparables de
datos epidemiológicos, como \textit{Our World in Data}
\cite{owid-coronavirus}.

\newpage

\section{Objetivos}

\subsection{Objetivo General}

Desarrollar un modelo de simulación Monte Carlo, basado en el procedimiento
propuesto por Xie \cite{Xie2020}, para estimar la cantidad esperada de casos
nuevos de COVID-19 en distintas zonas geográficas.

\subsection{Objetivos Específicos}

\begin{itemize}
  \item Adaptar el procedimiento Monte Carlo de Xie \cite{Xie2020} a un
        conjunto de datos temporales de COVID-19 organizado por zona geográfica.
  \item Recopilar y estructurar series de tiempo con variables relevantes,
        tales como casos nuevos, personas infectadas, recuperadas, casos
        activos y población de referencia por zona, utilizando fuentes
        comparables \cite{owid-coronavirus}.
  \item Modelar la probabilidad diaria de contagio como un componente aleatorio
        dependiente del tiempo y calibrado con la evolución observada de los
        contagios en cada zona.
  \item Comparar las trayectorias simuladas con los datos reales observados
        para identificar discrepancias, patrones recurrentes y posibles ajustes
        del modelo.
  \item Estimar, a partir de múltiples simulaciones, el valor esperado de
        casos nuevos por periodo en cada zona analizada.
\end{itemize}

\newpage

\section{Marco Teórico}

La simulación de epidemias puede abordarse desde modelos deterministas, como
los de tipo SIR, o desde modelos estocásticos que incorporan explícitamente la
incertidumbre del proceso de contagio.
Los modelos estocásticos consideran que la transmisión
ocurre a través de eventos aleatorios cuya realización concreta puede variar
incluso bajo condiciones iniciales similares. Esta segunda perspectiva resulta
especialmente útil cuando se desea representar fluctuaciones diarias,
heterogeneidad espacial y cambios abruptos en el comportamiento social.

En el trabajo de Xie \cite{Xie2020},
la propuesta modela la propagación de COVID-19 como un proceso de puntos en el
que cada persona infecciosa puede generar un número aleatorio de nuevos
contagios. La cantidad de infecciones secundarias se representa mediante una
distribución de Poisson de una tasa de infección \(R_t\),
mientras que el momento en que ocurre cada nuevo contagio se modela mediante
una distribución binomial negativa con media \(\mu_T\). Esta construcción
permite simular tanto la cantidad como el tiempo de aparición de los nuevos
casos.

En el código suplementario del artículo, \(R_t\) se usa como la tasa de
infección, mientras que la cantidad de casos nuevos diarios se registra por
separado en el vector \texttt{kk}. Para mantener esa misma lógica y evitar
confusiones, en este anteproyecto se definirá \(K_{z,t}\) como la variable
aleatoria que representa la cantidad de casos nuevos en la zona \(z\) durante
el día \(t\), e \(I_{z,t}\) como el número de casos infecciosos presentes en
esa misma zona y momento. Siguiendo la idea básica del artículo de
Xie \cite{Xie2020}, se tomará \(R_t\) como la tasa de infección diaria; al
trabajar por zonas, esta notación se adaptará como \(R_{z,t}\). Así, una
formulación simple del modelo es
\[
  K_{z,t} \sim \operatorname{Poisson}(R_{z,t} I_{z,t}),
\]
por lo que el valor esperado de casos nuevos queda dado por
\[
  E[K_{z,t}] = R_{z,t} I_{z,t}.
\]
En consecuencia, el objetivo principal de la simulación será aproximar
\(E[K_{z,t}]\), es decir, la cantidad esperada de casos nuevos para cada zona y
cada periodo.

Un punto central del modelo es que \(R_t\) no se considera constante. La tasa
de transmisión cambia con el tiempo para reflejar variaciones en el contacto
social y el comportamiento de la población. En el artículo de referencia, la
selección de patrones para \(R_t\) se apoya en el ajuste a datos observados
\cite{Xie2020}. Esta idea es especialmente valiosa para el presente proyecto,
porque permite representar cambios en la intensidad de la transmisión sin
abandonar el carácter probabilístico del modelo.

Desde el punto de vista empírico, la simulación requiere series de tiempo
comparables. No basta con observar el número acumulado de casos; es necesario
considerar variables que cambian día con día, tales como el número de personas
activamente infectadas, recuperadas, la población potencialmente expuesta y los
casos nuevos reportados en cada zona. Para documentar estos factores, resultan
particularmente útiles bases de datos como \textit{Our World in Data}
\cite{owid-coronavirus}, que reúne indicadores epidemiológicos comparables.

El proyecto también se apoya en la idea de validación por contraste entre
simulación y observación.
Si el modelo captura adecuadamente la tendencia general de contagios y la
variación de los casos nuevos en cada zona, entonces puede utilizarse para
obtener estimaciones esperadas a partir de muchas realizaciones del proceso
aleatorio. Si no lo hace, las discrepancias mismas aportan información sobre
variables omitidas o supuestos que deben refinarse.

\newpage

\section{Metodología}

\subsection{Recopilación y Organización de Datos}

La base empírica utilizada corresponde a los archivos locales del Oxford
COVID-19 Government Response Tracker, disponibles en el directorio
\texttt{data/} para Australia, Brasil, Canadá, China, Reino Unido, India y
Estados Unidos \cite{Hale2021}. Para evitar cargar columnas de notas y
variables no utilizadas, el procesamiento lee únicamente
\texttt{CountryName}, \texttt{CountryCode}, \texttt{Jurisdiction},
\texttt{Date} y \texttt{ConfirmedCases}. En cada archivo se conservaron solo
las observaciones nacionales, identificadas por
\texttt{Jurisdiction = NAT\_TOTAL}.

La variable \texttt{ConfirmedCases} es acumulada, por lo que los casos nuevos
observados se obtuvieron mediante diferencias diarias no negativas. Como los
archivos no contienen directamente una columna \(R_t\), se estimó un
\(R_t\) efectivo a partir del cociente entre el promedio móvil de siete días de
casos nuevos y el mismo promedio desplazado cinco días. Este desplazamiento
representa el intervalo generacional medio usado en la simulación. Para reducir
la inestabilidad causada por divisiones sobre bases muy pequeñas, el valor de
\(R_t\) se limitó al intervalo \([0,2]\). La simulación de cada país inició en
el primer día con al menos 100 casos nuevos reportados.

\subsection{Construcción del Modelo Estocástico}

Tomando como referencia el trabajo de Xie \cite{Xie2020}, se implementó un
modelo en el que los casos nuevos simulados de un día actúan como fuente de
contagios secundarios futuros. Para un día \(t\), si la simulación tiene
\(I_t\) infectados nuevos, la cantidad total de contagios secundarios se dibuja
como
\[
  S_t \sim \operatorname{Poisson}(I_t R_t).
\]
Esta forma es equivalente a dibujar una variable
\(\operatorname{Poisson}(R_t)\) por cada infectado y sumar los resultados, pero
es más eficiente computacionalmente. A cada contagio secundario se le asigna
un retraso \(D\) en días mediante una distribución binomial negativa con media
5 y dispersión 2, más un día adicional para evitar contagios en el mismo día.

Bajo esta notación, la simulación Monte Carlo produce realizaciones
\[
  K_{z,t}^{(1)}, K_{z,t}^{(2)}, \dots, K_{z,t}^{(N)},
\]
y el valor de expectación de casos nuevos
\[
  E[K_{z,t}] \approx
  \sum_{k=1}^{N} \frac{K_{z,t}^{(k)}}{N}.
\]

El procedimiento implementado fue el siguiente:

\begin{itemize}
  \item Para cada país se construyó una serie diaria con fecha, casos nuevos
        observados y \(R_t\) estimado.
  \item Se usaron los casos nuevos del primer día de la ventana como semilla
        inicial.
  \item En cada día observado, los casos simulados generaron contagios
        secundarios con distribución de Poisson.
  \item Los contagios secundarios fueron agregados a días futuros mediante la
        distribución de retrasos.
  \item Se extendió el horizonte 30 días más allá de la última fecha observada.
  \item El proceso se repitió 1000 veces por país y se calcularon percentiles
        5, 25, 50, 75 y 95 de casos diarios y acumulados.
\end{itemize}

\section{Resultados}

La Tabla \ref{tab:simulation-results} resume los resultados principales de las
1000 simulaciones por país. Las fechas iniciales varían porque cada serie se
recortó al primer día con al menos 100 casos nuevos. En todos los casos la
fecha final observada fue el 31 de diciembre de 2022 y el horizonte simulado se
extendió 30 días adicionales.

\begin{table}[htbp]
  \centering
  \small
  \begin{tabular}{lrrrr}
    \hline
    País & \(R_t\) medio & Pico mediano & Fecha modal del pico
    & Acumulado mediano \\
    \hline
    Australia & 1.069 & 237548 & 2022-01-19 & 37015378 \\
    Brasil & 1.059 & 79706 & 2022-02-05 & 18278252 \\
    Canadá & 1.039 & 21632 & 2022-01-12 & 3057878 \\
    China & 1.080 & 632952 & 2022-12-06 & 25822344 \\
    Reino Unido & 1.047 & 92554 & 2022-02-07 & 15543702 \\
    India & 1.032 & 110546 & 2021-05-12 & 15048894 \\
    Estados Unidos & 1.044 & 74788 & 2022-01-17 & 11410956 \\
    \hline
  \end{tabular}
  \caption{Resumen de 1000 simulaciones Monte Carlo por país.}
  \label{tab:simulation-results}
\end{table}

Los resultados muestran que el modelo reproduce de manera cualitativa algunos
periodos de crecimiento observados, pero no debe interpretarse como una
reconstrucción exacta de los casos reales. La diferencia principal es que la
simulación se alimenta de su propia trayectoria después del día inicial; por
ello, pequeñas diferencias tempranas pueden amplificarse durante ventanas de
más de mil días. Esto es visible en países como Australia y China, donde el
acumulado mediano simulado supera ampliamente los casos observados en la serie
procesada.

En contraste, Canadá, Reino Unido, India y Estados Unidos presentan acumulados
simulados menores que los observados. Esto sugiere que una sola cadena de
transmisión inicial, aun con \(R_t\) variable, no captura completamente la
entrada repetida de brotes, cambios de reporte, importaciones de casos ni
heterogeneidad regional. El resultado es útil como ejercicio de Monte Carlo
basado en \(R_t\), pero una versión predictiva debería incorporar reinyección
de casos observados, calibración por ventanas y posiblemente parámetros por
subregión.

\begin{figure}[htbp]
  \centering
  \includegraphics[width=0.9\textwidth]{generated/figures/USA_daily_cases.png}
  \caption{Casos nuevos observados y simulación Monte Carlo para Estados
  Unidos.}
  \label{fig:usa-daily}
\end{figure}

\begin{figure}[htbp]
  \centering
  \includegraphics[width=0.9\textwidth]{generated/figures/BRA_daily_cases.png}
  \caption{Casos nuevos observados y simulación Monte Carlo para Brasil.}
  \label{fig:brazil-daily}
\end{figure}

\newpage

\printbibliography

\end{document}
